\section{الفصل~\ref{sec:nora}. \nt{NORA}}

بنية نظام \nt{NORA}، ووضعاه، وصلته بالحدقة (لا عبر \nt{NORA} وحده)، وارتفاع نسبة الإشارة إلى الخلفية، وحرف U المقلوب في السلوك، كلها مقيسة مباشرة؛ والكسب الضربي $g$ نموذج؛ وأن الكسب يبدّل وضع الاختيار ويعمل مقلوبًا لدرجة الحرارة استنتاجان في الفصل؛ وفائدة ضبط الكسب حساب بنموذج الكتاب؛ أما ”\nt{NORA} مقبض الاستكشاف“ و”\nt{NORA} مقاطعة“ ففرضيتان.

{\footnotesize
\setlength{\tabcolsep}{3pt}\renewcommand{\arraystretch}{1.15}
\begin{longtable}{>{\raggedright}p{0.5cm}>{\raggedright}p{4.0cm}>{\raggedright}p{5.6cm}>{\raggedright}p{2.3cm}>{\raggedright\arraybackslash}p{1.7cm}}
\toprule
م & القضية & التجربة وما قيس & المصدر & الحالة \\
\midrule\endfirsthead
\toprule
م & القضية & التجربة وما قيس & المصدر & الحالة \\
\midrule\endhead
1 & عصبونات \nt{NORA} عند الإنسان بضع عشرات الآلاف، كلها تقريبًا في تجمّع واحد & مقاطع متسلسلة من جذع الدماغ البشري، صبغ بالتيروزين هيدروكسيلاز: $53\,900$ خلية في الموضع الأزرق، و$6\,260$ أخرى في النواة الفرعية المجاورة. & \cite{Baker1989LC} & مُقاس \\
2 & المخارج تتفرق في الدماغ كله تقريبًا، لكن أجزاء التجمع تخدم جزئيًا مناطق مختلفة & وسم جيني فيروسي عند الفئران: عصبونات \nt{NORA} في الموضع الأزرق تتلقى مداخل من مناطق واسعة في الدماغ وترسل مخارج إلى مناطق واسعة في الدماغ والنخاع الشوكي. وعند الجرذان: المجموعة الذاهبة إلى اللوزة تقوّي تعلّم الخوف، ومجموعة مستقلة ذاهبة إلى القشرة الجبهية الأمامية تقوّي إطفاءه. & \cite{Schwarz2015LC, Uematsu2017LC, Poe2020} & مُقاس \\
3 & الوضع الخلفي: نبضات منتظمة بتردد من أجزاء النبضة إلى بضع نبضات في الثانية، يتغير مع اليقظة & جرذان، تسجيل عصبونات الموضع الأزرق في دورة النوم واليقظة: التردد أعلى ما يكون في اليقظة، وأدنى في النوم البطيء، وقريب من الصفر في النوم السريع؛ وتغيرات التردد تسبق تغيّر الحالة. & \cite{AstonJonesBloom1981} & مُقاس \\
4 & الوضع الطوري: رشقة قصيرة لحدث مهم للمهمة، نحو لحظة القرار & قرود، مهمة البحث عن هدف نادر: كل عصبونات الموضع الأزرق تستجيب برشقة للهدف وحده، بعده بـ$\sim90\ms$ وقبل الاستجابة باليد بـ$\sim200\ms$؛ وعند الأداء السيئ تكون الرشقات أضعف. وفي مهمة الاختيار بين اثنين ترتبط الرشقة بالاستجابة أدق من ارتباطها بالمنبّه، وتظهر مع الاستجابات الصحيحة والخاطئة معًا. & \cite{AstonJones1994vigilance, Clayton2004LC} & مُقاس \\
5 & قطر الحدقة نقطة اختبار غير دقيقة لـ\nt{NORA}: فهو يتبع \nt{NORA} و\nt{ACET} ومناطق أخرى & قرود: نبضات الموضع الأزرق، حتى النبضات المفردة لخلية واحدة، تتنبأ باتساع الحدقة؛ وتوجد صلة مشابهة لكن متأخرة في الأكيمات وفي القشرة الحزامية. فئران: تقلبات الحدقة تتبع نشاط المخارج النورأدرينالية والكولينية في القشرة معًا. & \cite{Joshi2016, Reimer2016} & مُقاس \\
6 & النورأدرينالين يكبت النشاط الخلفي أكثر من المستثار؛ والكسب الضربي~\eqref{eq:gain} نموذج ينقل هذا الأثر & قرود يقظة، القشرة السمعية، رحلان أيوني للنورأدرينالين: يُكبت النشاط الخلفي للعصبونات أكثر من استجابتها لأصوات بني جنسها، فترتفع نسبة الإشارة إلى الضجيج. الدالة السينية الضربية~\eqref{eq:gain} مأخوذة من نموذج شبكي؛ والضرب الحرفي للانحدار غير مقيس. & \cite{Foote1975NE, ServanSchreiber1990} & مُقاس؛ $g$ نموذج \\
7 & الكسب يرفع $d'$ ضد الضجيج بعد المضخّم وحده~\eqref{eq:gainsnr} (القضية~\ref{prop:gainnoise}) & اشتقاق في الفصل؛ وفي النموذج الشبكي يحسّن الكسب كشف الإشارة. لا توجد تجربة تفصل الضجيج قبل المضخّم عن الضجيج بعده وتختبر هذا الفرق. & اشتقاق في الفصل؛ \cite{ServanSchreiber1990} & نظرية \\
8 & الحد $g\beta=1$ ينقل شبكة الاختيار من ”تتروّى“ إلى ”قررت“~\eqref{eq:wtagain} (القضية~\ref{prop:gainwta}، الشكل~\ref{fig:pitchfork}) & تحليل خطي لمجموعتين تثبّط إحداهما الأخرى؛ والاشتقاق في الفصل. لا توجد قياسات مباشرة لهذا الحد في الدماغ. & اشتقاق في الفصل & نظرية \\
9 & معامل الكسب مقلوب درجة حرارة الاختيار~\eqref{eq:softmax} & مع ضجيج لوجستي بعد المضخّم يكون احتمال الاختيار دالة \foreignlanguage{english}{softmax} مع $T=\sigma/g$؛ والاشتقاق في الفصل. & اشتقاق في الفصل & نظرية \\
10 & \nt{NORA} يحدد كم نستكشف: النشاط الخلفي المرتفع $g$ فعّال صغير (استكشاف)، والرشقة الطورية لحظة القرار $g$ كبير (قرار) & بشر: قبل الاختيار العشوائي تكون الحدقة القاعدية أوسع منها قبل اختيار أفضل خيار معروف؛ والأشخاص ذوو الحدقة الأوسع أميل إلى الاستكشاف. جرذان: الانتقال إلى وضع الاختيار العشوائي يتحكم فيه دخل \nt{NORA} إلى القشرة الحزامية الأمامية. أما أن الرشقة ترفع $g$ الفعّال فلم يُقَس مباشرة. & \cite{JepmaNieuwenhuis2011, Tervo2014stochastic, AstonJones2005} & فرضية \\
11 & المكافأة تعتمد على $g$ على شكل حرف U مقلوب؛ وأفضل $g$ أصغر في العالم السريع التغير (القضية~\ref{prop:explore}، الشكل~\ref{fig:invertedU}) & \texttt{models/nora\_explore.py}، $60$ محاكاة في كل منها $4000$ محاولة. تغيّر كل $1000$ محاولة: أفضل $g=16$، والنسبة $0.976$؛ وكل $20$: $g=8$، والنسبة $0.886$؛ وعند $g=0.5$: $0.77$. إعادة الحساب طابقت الفصل. & \texttt{models/\allowbreak nora\_\allowbreak explore.py} & نموذج \\
12 & حرف U المقلوب في السلوك: أفضل أداء عند نشاط خلفي معتدل مع رشقات واضحة & قرود، المهمة نفسها في السطر~4: في فترات الأداء الجيد يكون التردد الخلفي معتدلًا والرشقات للهدف واضحة؛ وعند التردد الخلفي المرتفع تغيب الرشقات وتكثر الاستجابات الخاطئة. وتغيرات النشاط الخلفي والمستثار تتبع تقلبات جودة الأداء. & \cite{Usher1999LC, AstonJones2005} & مُقاس \\
13 & \nt{NORA} يخبر بعدم اليقين غير المتوقع، و\nt{ACET} بالمتوقع & نظرية استدلال بايزي بنوعين من عدم اليقين؛ ولا توجد في الأعمال المستشهد بها تجربة مباشرة تفصل هاتين الإشارتين بحسب النواقل. & \cite{YuDayan2005} & فرضية \\
14 & بعد تغيّر حاد يتعلم الناس أسرع، وتتسع الحدقة & بشر، مهمة تنبؤ مع تغيرات مفاجئة: التغيرات القصيرة في الحدقة تتبع تقدير احتمال التغيّر، وتتنبأ مع الحدقة القاعدية بمدى تغيير المعطيات الجديدة للتقدير (معدل التعلّم)؛ وتغيير الحدقة بما لا يتصل بالضوء والمهمة يغيّر هذا المعدل. أما أن الحدقة تبيّن هنا \nt{NORA} تحديدًا فاستنتاج من المعطيات غير المباشرة في السطر~5. & \cite{Nassar2012} & مُقاس \\
15 & الرشقة الطورية تعمل مقاطعةً: تعيد الشبكة ضبط حالتها & لا توجد تجربة مباشرة تعيد فيها رشقة \nt{NORA} ضبط حالة الشبكة؛ والمقيس الرشقات عند الأحداث المهمة فقط (السطر~4). & \cite{BouretSara2005} & فرضية \\
16 & ضبط $g$ أو معدل التعلّم بحسب المفاجأة لا يكاد يفضل أفضل الإعدادات الثابتة & \texttt{models/nora\_explore.py}، عالم بحقب من $1000$ محاولة (تغيّر كل $1000$ وكل $20$). أفضل $g$ ثابت $=8$: $0.925$؛ وضبط $g$: حتى $0.927$؛ وضبط معدل التعلّم: حتى $0.926$؛ ومعدل ثابت $0.4$: $0.928$. إعادة الحساب طابقت الفصل. & \texttt{models/\allowbreak nora\_\allowbreak explore.py} & نموذج \\
\bottomrule
\end{longtable}}
